\section{Discussion}

\paragraph{The inverted-U as explore--commit dynamics.}
The inverted-U positional pattern admits a natural interpretation in terms of \emph{explore--commit} dynamics. Early in the reasoning chain, the model sets up the problem and has not yet engaged in substantive deliberation, producing few uncertainty markers. In mid-chain positions, the model actively evaluates alternatives, entertains competing hypotheses, and weighs evidence---activities that naturally produce hedging language. As the chain progresses toward a conclusion, the model commits to an answer and the need for hedging diminishes. This trajectory parallels deliberative processes studied in human cognition and decision-making, though we caution against strong cognitive analogies: the observed patterns are linguistic, and their relationship to underlying model computations remains unclear.

\paragraph{Confidence filtering and overconfidence.}
The finding that over half of reasoning-stage uncertainty is suppressed before the final response has implications for the well-documented overconfidence of LLMs \citep{mielke2022reducing, zhou2024relying}. If models are generating uncertainty expressions during reasoning but systematically removing them from responses, this suggests that the ``overconfidence problem'' is partly a \emph{presentation} issue rather than a purely epistemic one. The model's internal process contains hedging; the output does not. This distinction could inform approaches to calibration: rather than training models to express more uncertainty from scratch, it may be productive to reduce the filtering that occurs between reasoning and response.

The category-specific filtering patterns are informative. Epistemic hedges (``I think,'' ``perhaps'') are the most aggressively filtered, consistent with the view that these markers are perceived as inappropriate for authoritative-sounding responses. Approximators, by contrast, slightly increase---they serve a useful communicative function (``approximately 3 hours'') that survives or is even added during response generation.

\paragraph{Cross-model consistency.}
The positional inverted-U pattern is remarkably consistent across model families, despite large differences in overall uncertainty magnitude. This suggests that the pattern reflects a general property of autoregressive reasoning-trace generation rather than a family-specific training artifact. Families with high uncertainty rates (Minimax, Qwen) may simply have been trained with more hedging in their data or tuned to express more uncertainty, but the \emph{dynamics}---how uncertainty waxes and wanes across the chain---appear universal.

\paragraph{Limitations.}
Several limitations should be noted. First, our uncertainty detection is purely linguistic, based on surface markers. A sentence may express uncertainty without using any of our lexical patterns, and some matches may be false positives in context. Our validation against an LLM classifier (Cohen's $\kappa = 0.67$) suggests reasonable but imperfect coverage. Second, the dataset is observational: reasoning trace availability varies dramatically across model families, and the interactions are not controlled for topic, difficulty, or user intent. Third, we analyze only English-language markers; multilingual reasoning traces may exhibit different patterns. Fourth, the \textsc{real-slop} corpus reflects a specific time period and user population, limiting generalizability. Finally, we make no causal claims: the positional patterns and filtering effects are descriptive, and alternative explanations (e.g., positional effects driven by sentence length or syntactic complexity) cannot be ruled out.
